\documentclass{article}
\usepackage[UTF8]{ctex}
\usepackage{amsmath,amssymb,amsfonts}
\usepackage{enumitem}
\usepackage{geometry}
\geometry{a4paper, margin=1in}
% === 错题本排版优化（不参与题目标识） ===
\setlength{\parindent}{0pt}
\setlength{\parskip}{0.35em plus 0.10em minus 0.08em}
\setlength{\emergencystretch}{3em}
\setlist{itemsep=0.25em, topsep=0.40em, parsep=0pt}
\allowdisplaybreaks[2]
% === 错题本排版优化结束 ===
\begin{document}
\section*{南开大学2025年数学分析真题}
\section*{56．南开大学2025年数学分析试卷}

\subsection{题目1}
1．（25 分）求极限

\[
\lim _{x \rightarrow 0} \frac{\sin (\ln (1+x))-\ln (1+\sin x)}{\sin (\sin x)-\ln (1+\ln (1+x))}
\]

\subsection{题目2}
2．（25 分）设函数 \(f(x)\) 在点 \(x_{0}\) 处连续，在 \(x_{0}\) 的去心邻域内二阶可导且二阶导函数有界．证明：\(f(x)\)在点 \(x_{0}\) 处左右导数都存在．

\subsection{题目3}
3．（25 分）设 \(f(x)\) 在 \([0,+\infty)\) 上连续，且 \(\lim _{x \rightarrow+\infty} f(x)=A, g(x)\) 在 \([0,1]\) 上可积，证明：

\[
\lim _{n \rightarrow \infty} \int_{0}^{1} f(n x) g(x) \mathrm{d} x=A \int_{0}^{1} g(x) \mathrm{d} x
\]

\subsection{题目4}
4．（20 分）设 \(f:(a, b) \rightarrow(a, b)\) 满足对任意的 \(x, y \in(a, b)\) ，当 \(x \neq y\) 时，有

\[
|f(x)-f(y)|<|x-y| .
\]

任取 \(x_{1} \in(a, b)\) ，令 \(x_{n+1}=f\left(x_{n}\right), n=1,2, \cdots\) ，证明：数列 \(\left\{x_{n}\right\}_{n=1}^{\infty}\) 收敛。

\subsection{题目5}
5．（20 分）设定义在 \(\mathbb{R}\) 上的函数 \(f(x)\) 在 \(x=0\) 处连续，对任意满足 \(a_{n}<0<b_{n}\) 且 \(\lim _{n \rightarrow \infty} a_{n}=\) \(\lim _{n \rightarrow \infty} b_{n}=0\) 的点列 \(\left\{a_{n}\right\},\left\{b_{n}\right\}\) ，极限 \(\lim _{n \rightarrow \infty} \frac{f\left(b_{n}\right)-f\left(a_{n}\right)}{b_{n}-a_{n}}\) 都存在．证明：函数 \(f(x)\) 在 \(x=0\) 处可微．

\subsection{题目6}
6．（20 分）设 \(f(x)=\arctan (\sin x), x_{1}=\frac{1}{2}\) ，令 \(x_{n+1}=f\left(x_{n}\right), n=1,2, \cdots\) ，讨论 \(\sum_{n=1}^{\infty} x_{n}^{2}\) 的敛散性。

\subsection{题目7}
7．（15 分）求重积分 \(\iint_{x^{2}+y^{2} \leq 1}\left(e^{y}+e^{-y}\right) \cos x \mathrm{~d} x \mathrm{~d} y\) ．
\end{document}
